% !TEX root = ../main.tex
% File: part1/chapters_reasoning/sec3_self_improvement.tex

\section{Tự học Suy luận: STaR, ReST\texorpdfstring{$^{EM}$}{-EM} và Quiet-STaR \newtag}
\label{sec:self_improvement}

Dữ liệu chuỗi suy luận do con người viết rất đắt và khó mở rộng. Một ý tưởng hấp dẫn: để mô hình \textbf{tự sinh ra lời giải}, \textbf{giữ lại những lời giải đúng}, rồi \textbf{học lại từ chính chúng}. Ba công trình dưới đây là các bậc thang dẫn tới học tăng cường cho suy luận.

\subsection{STaR: Tự khởi động suy luận bằng suy luận}
\label{ssec:star}
\textbf{STaR} (Self-Taught Reasoner) \cite{zelikman2022star} chỉ cần một vài ví dụ lời giải mẫu và một tập lớn câu hỏi có đáp án (không có lời giải). Vòng lặp:
\begin{enumerate}
	\item \textbf{Sinh lời giải:} Dùng vài ví dụ mẫu làm few-shot prompt, cho mô hình sinh lý giải (rationale) và đáp án cho mọi câu hỏi.
	\item \textbf{Lọc:} Chỉ giữ lại các lý giải dẫn tới \emph{đáp án đúng}.
	\item \textbf{Hợp lý hóa (rationalization):} Với các câu mô hình trả lời sai, đưa \emph{đáp án đúng làm gợi ý} vào prompt và yêu cầu mô hình viết lý giải dẫn tới đáp án đó. Nếu lý giải hợp lệ, thêm vào dữ liệu (như thể mô hình tự nghĩ ra, không kèm gợi ý). Bước này giúp mô hình học được từ những bài nó chưa giải được.
	\item \textbf{Tinh chỉnh:} Tinh chỉnh lại \emph{mô hình gốc ban đầu} trên toàn bộ lý giải đã lọc, rồi lặp lại từ bước 1 với mô hình mới.
\end{enumerate}
Trên CommonsenseQA, STaR với GPT-J 6B đạt kết quả tương đương việc tinh chỉnh một mô hình lớn hơn 30 lần. STaR có thể được hiểu như một xấp xỉ của \emph{policy gradient} với phần thưởng nhị phân (đúng/sai), trong đó việc lọc đóng vai trò loại bỏ các mẫu có phần thưởng bằng 0.

\subsection{ReST\texorpdfstring{$^{EM}$}{-EM}: Tự huấn luyện dưới góc nhìn Expectation-Maximization}
\label{ssec:restem}
\textbf{ReST$^{EM}$} \cite{singh2024restem} (Google DeepMind) đơn giản hóa và mở rộng quy mô ý tưởng trên, đồng thời đặt nó trên nền tảng lý thuyết EM cho học tăng cường với phần thưởng nhị phân:
\begin{itemize}
	\item \textbf{Bước E (Generate):} Lấy mẫu nhiều lời giải cho mỗi bài toán từ mô hình hiện tại; dùng phản hồi tự động (đáp án đúng cho toán, test case cho mã) làm phần thưởng nhị phân để lọc.
	\item \textbf{Bước M (Improve):} Tinh chỉnh mô hình để cực đại log-likelihood của các lời giải đã lọc (có trọng số theo phần thưởng). Giống STaR, mỗi vòng tinh chỉnh bắt đầu lại từ \emph{mô hình tiền huấn luyện gốc} thay vì mô hình của vòng trước, để giảm quá khớp.
\end{itemize}
Thí nghiệm với các mô hình PaLM 2 trên toán (MATH) và lập trình (APPS) cho thấy ReST$^{EM}$ \textbf{mở rộng tốt theo kích thước mô hình} và \textbf{vượt đáng kể việc tinh chỉnh chỉ trên dữ liệu người viết}. Chỉ vài vòng lặp là đủ -- nhiều vòng hơn có thể gây quá khớp trên tập huấn luyện. Thông điệp: với các tác vụ có thể tự động kiểm chứng, \emph{dữ liệu do mô hình tự sinh có thể tốt hơn dữ liệu con người}.

\subsection{Quiet-STaR: Suy nghĩ trước khi nói, ở mọi token}
\label{ssec:quiet_star}
STaR và ReST$^{EM}$ học suy luận cho các \emph{bài toán hỏi--đáp}. \textbf{Quiet-STaR} \cite{zelikman2024quietstar} đặt câu hỏi tổng quát hơn: văn bản thông thường cũng chứa nhiều suy luận ngầm “giữa các dòng” -- liệu mô hình có thể học cách tự suy nghĩ \emph{trước khi dự đoán mỗi token} trên văn bản internet bất kỳ?

\paragraph{Cơ chế.}
\begin{enumerate}
	\item \textbf{Nghĩ:} Tại \emph{mỗi} vị trí trong chuỗi, mô hình sinh một đoạn “suy nghĩ” ngắn, được bao bởi hai token đặc biệt học được \texttt{<|startofthought|>} và \texttt{<|endofthought|>}. Để làm việc này song song cho mọi vị trí, các tác giả đề xuất một thuật toán lấy mẫu song song theo token dựa trên mặt nạ attention đặc biệt.
	\item \textbf{Nói:} Một \textbf{đầu trộn (mixing head)} nhỏ học trọng số để nội suy giữa dự đoán \emph{có} suy nghĩ và dự đoán \emph{không có} suy nghĩ -- giúp giai đoạn đầu huấn luyện không bị các suy nghĩ còn vô nghĩa làm hỏng.
	\item \textbf{Học:} Dùng REINFORCE: một suy nghĩ được thưởng nếu nó làm \emph{tăng} log-likelihood của các token thật tiếp theo so với mức trung bình của các suy nghĩ khác tại cùng vị trí. Kỹ thuật teacher-forcing mở rộng cho phép đánh giá tác dụng của suy nghĩ lên nhiều token tương lai chứ không chỉ một token.
\end{enumerate}

\paragraph{Kết quả.} Sau khi tiếp tục tiền huấn luyện Mistral 7B trên văn bản internet bằng Quiet-STaR, độ chính xác \textbf{zero-shot} (không hề tinh chỉnh trên các tác vụ này) tăng từ 5.9\% lên 10.9\% trên GSM8K và từ 36.3\% lên 47.2\% trên CommonsenseQA. Các token khó dự đoán hưởng lợi nhiều nhất từ việc “suy nghĩ”.

\begin{tcolorbox}[title={Từ tự học đến học tăng cường}, colback=blue!5!white, colframe=blue!60!black, fonttitle=\bfseries]
Cả ba phương pháp đều là các dạng \textbf{học tăng cường với phần thưởng nhị phân}, được thực hiện “gián tiếp” qua lọc dữ liệu (STaR, ReST$^{EM}$) hoặc REINFORCE (Quiet-STaR). Điểm hạn chế chung: chúng \emph{chỉ học từ lời giải đúng} (hoặc chỉ theo tín hiệu tương đối yếu), bỏ phí thông tin từ các lời giải sai. Học tăng cường trực tuyến với advantage có dấu (mục tiếp theo) tận dụng được cả hai phía: đẩy xác suất của lời giải tốt lên \emph{và} kéo xác suất của lời giải tệ xuống.
\end{tcolorbox}
